\scanpage{37}
\uline{Lemme \deleted{\unsure{2}}3}\quad Soit $G$ schéma simple et propre sur $S$ \deleted{\unsure{loc. nh.|loc. noeth.}}\inserted{$=\operatorname{Spec}(A)$} \deleted{l'anneau $A$}\inserted{$A$ artinien local}, $G$ muni d'une section $e$. \deleted{Soit $S_0$ \unsure{$=\operatorname{Spec}(A/I)$|$=\operatorname{Spec}(A/\mathfrak J)$}} \deleted{avec $\unsure{I=0}$, soit $I$ nilpotent, $I^2=0$, $I$ libre sur $S_0$}\inserted{On suppose $\mathfrak m^{n+1}=0$}. On suppose que $G_{n-1}=G\times_S S_{n-1}$ est un schéma abélien ayant $e_0$ comme section unité. \deleted{Alors il Soit} Alors il existe une structure multiplicative unique sur $G$, induisant celle de $G_{n-1}$ et pour laquelle $e\times e$ soit \deleted{unité à gauche (resp. à droite)}\inserted{s'applique en $e$}. De plus, $G$ est un schéma abélien pour cette structure, de section unité $e$.

\deleted{L'obstruction} À cause de l'unicité, on peut supposer $S$ \deleted{\unsure{est|noeth.}} affine. L'obstruction $\unsure{\zeta|\xi}$ au \unsure{prolongmt} de l'application $G_0\times G_0\xrightarrow{\pi_0}G_0$ est alors dans
\[
 H^1(G_0\times G_0,\pi_0^*(\unsure{\underline{\mathcal V}|\underline{\mathcal T}}_{G_0/S_0})\otimes\unsure{\underline{\mathfrak m}^{n}|\underline{M}^{n}})=.
\]
\deleted{Ici \uline{Künneth} s'applique}
\[
H^1(G_0\times G_0,\underline O_{G_0\times G_0})\otimes_{\deleted{\unsure{K_0|k_0}}k}(\mathcal V\otimes_k\underline{\mathfrak m}^{n})).
\]
Künneth s'applique, il suffit de montrer que la restriction de $\zeta$ à $e\times G_0$ et $G_0\times e$ est nulle. Cela est vrai, car ces \deleted{\unsure{obstructions}}\inserted{restric\unsure{tions}} $,\ \in H^1(G_0,\underline O_{G_0})\otimes_k(\mathcal V\otimes_k\mathfrak m^n)$ sont resp. les obstructions à prolonger $g\longrightarrowgg$ et $g\longrightarrow ge$, qui sont les applications identiques. On peut donc prolonger, l'indétermination étant un
\[
H^0(G_0\times G_0,\underline O_{G_0\times G_0})\otimes_k(\mathcal V\otimes_k\mathfrak m^n).
\]

% slanted
\margin{lemme 2 auxiliaire} Plus précisément, si on a déjà $X_{\unsure{n-1}}\longrightarrow G$, l'obstruction à prolonger est dans $H^1(X_0,\underline O_{X_0})\otimes_k\mathcal V\otimes_k\mathfrak m^n$, et l'indétermination dans $H^0(X_0,\underline O_{X_0})\otimes_k\mathcal V\otimes_k\mathfrak m^n$.
% archive sheet number: 36
